\documentclass[11pt]{report}
\usepackage[margin=1in]{geometry}
\usepackage{xcolor}
\usepackage{listings}
\usepackage{longtable}
\usepackage{booktabs}
\usepackage{enumitem}
\usepackage{hyperref}
\hypersetup{colorlinks=true, linkcolor=blue!50!black, urlcolor=blue!50!black}

\definecolor{codebg}{RGB}{245,245,248}
\lstset{
  basicstyle=\ttfamily\small, backgroundcolor=\color{codebg},
  breaklines=true, frame=single, rulecolor=\color{gray!40},
  columns=fullflexible, keepspaces=true, showstringspaces=false
}
\newcommand{\code}[1]{\texttt{#1}}

\title{\textbf{Canabalt}\\[4pt]\large A didactic C++ port of the one-button runner ---
       build, play \& mechanics}
\author{Built on the Otacon engine (see \code{docs/otacon-engine.pdf})}
\date{}

\begin{document}
\maketitle
\tableofcontents

\chapter{About Canabalt}
\emph{Canabalt} is the original one-button infinite runner: a lone figure sprints
across the rooftops of a collapsing city, and the only control is \textbf{jump}.
Hold longer to jump higher; clear the gaps; survive as far as you can. There is no
winning --- only your distance in metres, and your best.

This is a teaching-oriented C++ port. The source of truth is the official
open-source release (\code{ericjohnson/canabalt-ios}), an Objective-C game on the
\code{flixel-ios} framework, kept under \code{reference/} and studied closely so
the port does \textbf{not diverge} from the original's feel. The game is written
entirely against the Otacon engine through its \code{IGame} interface; everything
game-specific --- code, assets and these docs --- lives under
\code{src/game/canabalt/}.

\section{Faithfulness}
The simulation reproduces flixel's exact integrator, velocity model, swept-AABB
collision and follow camera, plus Canabalt's exact player physics constants
(Chapter~\ref{ch:mech}). With the default float scalar the run is deterministic
and reproducible; the same seed yields the same city every time.

\chapter{Building and running}
Canabalt is the \code{canabalt} executable; it links the Otacon static library.

\section{macOS / Linux}
\begin{lstlisting}[language=bash]
cmake -S . -B build            # first configure fetches the window library
cmake --build build -j8
./build/bin/canabalt
\end{lstlisting}

\section{Windows}
\textbf{Visual Studio:} open the folder (it has \code{CMakeLists.txt}); VS
configures automatically. Pick the \code{canabalt.exe} target and Run.\\
\textbf{VS~Code:} with the \emph{CMake Tools} extension, select a kit, then Build
and Run. Or from a terminal: \code{cmake -S . -B build \&\& cmake --build build}.

\section{Vulkan note}
If built with the Vulkan backend on macOS, launch through \code{tools/run.sh}
(it sets the MoltenVK loader environment). Any backend plays identically.

\chapter{How to play}
You start on the \textbf{title screen}; press \textbf{jump} to begin. From then on
you run automatically and ever faster --- tap jump to hop a gap, hold it for a
higher leap. Falling into a gap, smacking a wall, or being caught by a hazard ends
the run.

\begin{itemize}[nosep]
  \item \textbf{Distance.} The metres-run counter sits top-right.
  \item \textbf{Game over.} On death a panel shows how far you ran and how you
        died (``\dots before falling to your death'', ``\dots smashing into a
        wall'', ``\dots turning into a fine mist''). Press \textbf{jump} or
        \code{R} to run again, or click \textbf{EXIT} (top-right) to return to the
        title.
  \item \textbf{High score.} Beat your best and a \textbf{NEW RECORD} badge
        appears; the best distance is saved to
        \code{src/game/canabalt/assets/highscore.dat} and persists between
        sessions.
\end{itemize}

\section{Controls}
\begin{longtable}{ll}
\toprule
Key / input & Action \\
\midrule
\endhead
Space / W / Up / Left-click & Jump (hold for a higher jump); start; retry \\
\code{R} & Retry the current run \\
Click \textbf{EXIT} (on game over) & Back to the title \\
\code{1} \dots \code{4} & Jump directly to that mode \\
\code{]} or \code{.} \;/\; \code{[} or \code{,} & Next / previous mode \\
\code{F9} & Mute / unmute all audio \\
\code{F10} & Cycle the music track (Run / Daring Escape / Mach Runner) \\
\code{M} & Show / hide the mouse cursor \\
\code{`} (backtick) & Hide / show ALL dev UI --- clean play \\
\code{Esc} & Quit \\
\bottomrule
\end{longtable}
The teaching toggles (\code{F1}--\code{F8}, \code{H}, \code{+/-}, \code{C}, \code{G},
\dots) are documented with the modes below and in the engine manual; the bottom
legend always shows the live state. Press \code{`} any time to hide all of it and
just play.

\chapter{The incremental modes}
The build is staged as four runnable ``modes'' that show the game assembled from a
single box up to the full runner. Cycle with \code{]}/\code{[}, or jump straight to
one with the number keys \code{1}--\code{4}.
\begin{enumerate}[nosep]
  \item \textbf{Static Box} --- a box on the ground. Geometry + camera only.
  \item \textbf{Gravity} --- a box falls and lands on a platform (the integrator
        and collision). Press Space to drop again; press \code{E} to burst
        particles at the cursor.
  \item \textbf{Run \& Jump} --- the real \code{Player} auto-runs on flat ground;
        hold to jump; the camera follows. Full original feel, no hazards.
  \item \textbf{Infinite Run} --- the real game: the ported level generator, the
        parallax city, hazards, sprites, sound and the title/score flow.
\end{enumerate}

\chapter{Game mechanics}\label{ch:mech}
\section{Player physics (from \code{Player.m})}
\begin{longtable}{ll}
\toprule
Quantity & Value \\
\midrule
\endhead
Hitbox & $16 \times 18$ (sprite $30\times30$, offset $6,12$) \\
Start velocity $x$ & 125 \\
Gravity (accel.\ $y$) & 1200 \\
Drag $x$ & 640 \\
Max velocity & $(1000, 360)$ \\
Accel.\ $x$ ramp & $60/36/24/12/4$ as speed rises (fast then tapering) \\
Jump limit & $v_x / (1000 \cdot 2.5)$, capped at $0.35$\,s \\
Short / full jump $v_y$ & $-0.65\cdot360$ / $-360$ \\
Fall death & $y > 484$ \\
Wall death & accel.\ $x \to 0$, $v_x \to 0$ \\
\bottomrule
\end{longtable}
The hold-to-jump window scales with speed: the faster you run, the higher you can
leap. A hard landing makes the runner stumble and roll.

\section{The level generator (\code{Sequence.cpp})}
A faithful port of the original \code{Sequence.m}. Exactly two ``sequences''
leapfrog: each owns one wide building, and when one scrolls off the left it
regenerates just past the other. Each new building's gap, minimum width and
vertical drop are derived from your current \code{velocity.x} and \code{jumpLimit}
--- so the level gets harder precisely as you get faster, exactly like the
original. The first building is a fixed launch hallway.

Special building archetypes are generated by type, each with its own collision
shape, art and behaviour:
\begin{itemize}[nosep]
  \item \textbf{Hallway} --- a tunnel you run through; its windows smash (glass +
        sound) as you pass, and doors line the back wall (the launch tunnel has none).
  \item \textbf{Collapse} --- the roof sinks the moment you step on it; you ride it
        down and must leap off before it drops away.
  \item \textbf{Bomb} --- a bomb whistles down and detonates on the roof; touching
        it is fatal.
  \item \textbf{Crane} --- a metal beam you cross like a platform (footsteps ring
        out as metal).
  \item \textbf{Billboard} --- a thin catwalk ledge under a damaged sign.
  \item \textbf{Leg} --- a giant mech leg stomps down and crushes the building below.
\end{itemize}
Rooftops also scatter obstacles (crates you stumble over) and pigeons that flush as
you approach. Roof and floor caps, wall and window styles all vary per building.

\section{Background \& foreground extras}
Distant Easter eggs ride the skyline --- a mothership, a dark tower, two marching
mech \emph{walkers} (which vent smoke), and a jet that periodically streaks past
and rattles the screen --- while a steel girder occasionally sweeps the foreground.
A startup screen-quake (the collapsing city) shakes the camera as a run begins.

\chapter{Presentation}
\section{Progressive sprite reveal (\code{F6} / \code{F8})}
The world can be peeled between pure collision geometry and full sprites one layer
at a time. \code{F6} cycles the level --- \textbf{geometry} (everything is boxes)
$\to$ \textbf{player} (the runner becomes its \code{player2.png} sprite) $\to$
\textbf{buildings} (rooftops tile a wall texture) $\to$ \textbf{background}
(parallax skyline) $\to$ \textbf{full sprites} (textured particles) --- and
\code{F8} jumps straight between all-geometry and all-sprites. \code{C} overlays
the collision boxes at any level, so you can see the geometry under the art.

\section{Sound \& music}
Sound effects (footsteps, jumps, window smashes, glass, the bomb/leg/collapse
hazards, the jet, pigeon flaps, the death thud) are decoded in-house from the
original \code{.caf} files and mixed live. Three looping music tracks are included;
\code{F9} mutes all audio and \code{F10} cycles the track. The title screen has its
own theme.

\chapter{Assets}
All game assets live under \code{src/game/canabalt/assets/} and are loaded through
the engine's in-house decoders --- nothing is baked into the binary:
\begin{itemize}[nosep]
  \item \code{images/} --- PNGs (player sheet, HUD glyphs, building/crane/billboard
        tiles, UI: title, game-over, paused, cursor). Decoded by the in-house PNG
        decoder; \code{images/raw/} holds the un-atlased source art.
  \item \code{sound/} --- sound effects as $22050$\,Hz mono $16$-bit LPCM
        \code{.caf}. \code{convert.sh} regenerates them from the \code{.mp3}
        sources via \code{afconvert}.
  \item \code{music/} --- the looping tracks, converted to \code{.caf} the same way.
\end{itemize}

\chapter{Teaching toggles (quick reference)}
These flip engine/game behaviour live so the effect is visible at once; the HUD's
second line shows the current state.
\begin{longtable}{ll}
\toprule
Key & Effect \\
\midrule
\endhead
\code{F1} & RNG: deterministic (fixed-seed LCG) vs.\ entropy-seeded; regenerates the city \\
\code{F2}/\code{F3}/\code{F4} & Startup screen-quake on-off / weaker / stronger \\
\code{F5} & Variable vs.\ fixed simulation timestep \\
\code{F6}/\code{F8} & Sprite-reveal level / all-geometry$\leftrightarrow$all-sprites \\
\code{F7} & Particles on / off \\
\code{F9}/\code{F10} & Sound mute / music track \\
\code{H} & Perf + memory overlay; hover an entity to inspect it \\
\code{+}/\code{-} & Simulation time scale (slow-mo / fast-forward) \\
\code{C}/\code{X}/\code{Z}/\code{G} & Colliders / parallax / wireframe / grid overlays \\
\code{P}/\code{O} & Pause / single-step \\
\code{F12} & Save a PNG screenshot \\
\bottomrule
\end{longtable}

\end{document}
